\begin{sectionbox}
    \subsection{Bewegungen}

    \subsubsection{Freier Fall}
    $h=\frac{1}{2}gt^2 \qquad v_E=\sqrt{2gh_0}$\\
    $\Delta t = \sqrt{2h/g}$

    \subsubsection{Senkrechter Wurf}
    $h=v_0t-\frac{1}{2}gt^2+h_0$\\
    $v_M=\sqrt{\frac{1}{2}v_0} \qquad v^2 = v_0^2-2g(h_0-h)$\\
    $v=v_0-g  t$\\
    $t_{h,\ir max}=\frac{v_0}{g} \qquad h_{\ir max}=\frac{v_0^2}{2g}$

    \subsubsection{Waagrechter Wurf}
    $t_{\ir Fall}=\sqrt{\frac{2h_0}{g}}$ \\
    $y(t)=h_0-\frac{1}{2}gt^2 \qquad x(t)=v_0  t$\\
    $v_y=gt \qquad v_x = v_0$\\
    $y(x)=h_0-\frac{g  x^2}{2v_0^2}$

    \subsubsection{Schräger Wurf}
    $v_x=v_0\cos \alpha \qquad v_y=v_0\sin \alpha -gt$\\
    $v(\Delta h)=\sqrt{v_0^2-2g\Delta h}$\\
    $x=v_0t\cos \alpha \qquad y=v_0t \sin \alpha - \frac{1}{2} gt^2 +h_0$\\
    $x_{\ir max}=\frac{v_0^2 \sin{(2\alpha)}}{g}$\\
    $t_{\ir Wurf}=\frac{2v_0\sin \alpha}{g}$ \quad
    $t_{h,\ir max}=\frac{1}{2}t_{\ir Flug}$ \quad
    $h_{\ir max}=\frac{v_0^2 \sin^2 \alpha}{2g}+h_0$ \\

    % \subsubsection{Schiefe Ebene}
    % $F_H=F_G \sin \alpha \qquad F_N=F_G\sin\alpha$ \\
    % $F_{HR}=\mu_{HR}F_N \qquad F_{GL} = \mu_{GL}F_N$ \\

    \begin{cookbox}{Newton'sche Axiome}
        \item Masse ist träge. Jeder Körper behält seinen Bewegungszustand bei.
        \item $\vec F=m\vec a=\dot{\vec p}$
        \item $\vec F_{12}=-\vec F_{21}$
    \end{cookbox}
\end{sectionbox}


\begin{sectionbox}
    \subsection{Kraft, Energie, Impuls}
    \subsubsection{Kraft}
    \begin{tablebox}{ll}
        Federkraft & $\vec F_F=-kx$ \\ & $x:$ Auslenkung vom GGP \\
        Gewichtskraft & $F_G=mg$ \\
        & $F_N=mg\cos \alpha$ \\
        & $F_H=mg\sin \alpha$ \\
        Zentripetalkraft & $F_Z=ma_{Z}=m \frac{v^2}{r}=m\omega^2r=p\omega$ \\
        & Zeigt in Richtung Mittelpunkt des Kreises von \\ & einem sich nicht rotierenden Beobachter \\
        Zentrifugalkraft & $F_{\it flieh}=-F_Z$ \\
        & Zeigt vom Mittelpunkt des Kreises weg. Der Beobachter \\ & rotiert sich hier mit ($\rightarrow$ Scheinkraft) \\
        Reibungskraft & $F_R=\mu F_{(N)}=\mu m g   \cos \alpha$
    \end{tablebox}
    \subsubsection{Energie}
    \begin{tabular}{c|c}
        $E_{\ir kin}=\frac{1}{2}mv^2=\frac{p^2}{2m}$ & $E_{\ir pot}=mgh$                  \\
        $E_{F}=\frac{1}{2}kx^2$                      & $E_{\ir rot}=\frac{1}{2}I\omega^2$
    \end{tabular}

    \subsubsection{Arbeit}
    $W=\Delta E=\vec F \vec s=|\vec F| \, |\vec s|\cos\alpha$
    \subsubsection{Leistung}
    $P=
        \underbrace{\frac{\Delta W}{\Delta t}}_{\text{mittlere Leistung}}
        = \, \underbrace{\frac{F  \Delta s}{\Delta t}=F  v}_{\text{momentane Leistung}}$\\
    $\eta=\frac{E_{\it nutz}}{E_{\ir zu}}=\frac{P}{P}$
    \subsubsection{Impuls}
    $p_{\ir vorher}=p_{\ir nachher} \qquad p=mv$
    \begin{tablebox}{ll}
        Mittlere Kraft     & $\overline{F}=\frac{\Delta p}{\Delta t}$\\
        Elastischer Stoß   & $m_1v_1+m_2v_2=m_1u_1+m_2u_2$\\
        & $E_{\ir kin,ges}=$ const. \\
        Unelastischer Stoß & $m_1v_1+m_2v_2=(m_1+m_2)u$ \\
        & $E_{\ir kin,ges}\ne$ const.
    \end{tablebox}
    \subsubsection{Raketengleichung}
    TODO
\end{sectionbox}


\begin{sectionbox}
    \subsection{Rotation}
    \begin{tabular}{l|ll}
        Kinematik       & Rotation                & Einheit[Rot.]            \\ \hline
        Masse $m$       & Trägheitsmoment $I$/$J$ & $[I]=kgm^2$              \\
        Impuls $\vec p$ & Drehimpuls $\vec L$     & $[L]=\frac{kgm^2}{s}=Js$ \\
        Kraft $\vec F$  & Drehmoment $\vec M$     & $[M] = Nm$
    \end{tabular}

    \subsubsection{Bewegungsgleichungen}
    \begin{tabular}{ll}
        Winkel                & $\varphi(t)=\omega t$ bei $\omega=\text{const.}, \alpha = 0$                 \\
                              & $\varphi(t)=\frac{1}{2} \alpha t^2$ bei $\alpha=\text{const.}$               \\
        Winkelgeschwindigkeit & $\omega = \frac{2\pi}{T}=\frac{\Delta \varphi}{\Delta t}$                    \\
                              & $\vec \omega = \frac{\vec r \times \vec v}{r^2}$ (in Richtung der Rot.achse) \\
        Winkelbeschleunigung  & $\alpha = \frac{\Delta \omega}{\Delta t}$                                    \\
        Bahngeschwindigkeit   & $\vec v = \vec \omega \times \vec r$
    \end{tabular}

    \subsubsection{Trägheitsmoment $I$ oder $J$}
    \begin{tablebox}{lll}
        Zylindermantel & $I = mr^2$ & Drehachse ist Körperachse\\
        Hohlzylinder & $I = \frac{1}{2}m(r_a^2+r_i^2)$ & Drehachse ist Körperachse\\
        Massive Kugel & $I = \frac{2}{5}mr^2$  & Drehachse durch Mittelpunkt\\
    \end{tablebox}
    TODO: Satz von Steiner

    \subsubsection{Drehimpuls}
    $\vec L=\vec r\times \vec p = \vec r \times m\vec v = I \vec \omega$ \\
    Immer im Bezug auf Ursprung von $\vec r$ mit $\vec r$: Abstand von der Drehachse

    \subsubsection{Drehmoment (analog zu Kraft)}
    $\vec M=\vec r\times \vec F=|\vec F| \, |\vec r|  \sin \alpha$\\
    $\vec M=I  \dot{\vec{\omega}}=I\vec \alpha=\frac{\partial \vec L}{\partial t}$
\end{sectionbox}

\begin{sectionbox}
    \subsection{Corioliskraft}
    \begin{emphbox}
        $\vec{F_c}=2m(\vec v\times \vec \omega)= 2m v w   \sin{(v\angle \omega)}$
    \end{emphbox}
    $\vec v$: Bewegungsrichtung, $\vec w$: (Richtung der) Rotationsachse, senkrecht zur Rotationsrichtung (auf der Erde nach Norden)\\
    $a_c=-2(\omega \times v)=2v\omega$\\
    $\Delta x= \frac{2}{3}\cos (\beta) \omega h \sqrt{\frac{2h}{g}}$\\
    $\Delta y = -\omega^2R\cos (\beta) \sin (\beta) \frac{h}{g}$\\
    %\includegraphics[width=0.25\columnwidth]{coriolis.pdf}
    $R_E=\SI{6370}{\kilo \meter} \qquad r=R_E \cos \varphi \qquad \omega_E=\frac{2\pi}{86400\si{\second}}$\\
    $F_C$: Auf der Nordhalbkugel nach rechts. Auf der Südhalbkugel nach links.
\end{sectionbox}

\begin{sectionbox}
    \subsection{Planetenbewegungen}
    \subsubsection{Gravitation}
    $E_{\ir pot}=-G\frac{m_E  m}{r} \qquad r>r_E$\\
    $\vec{F}_G=G\frac{m_1m_2}{r^2}  \vec{e_r}$\\
    $g=G  \frac{m}{r^2}$\\
    $W_G=Gm_1m_2(\frac{1}{r_1}-\frac{1}{r_2})$\\
    $F_G \, || \, r \rightarrow M=r\times F=0$\\

    Satellit-Erde: $\frac{1}{2}m_Sv^2=G\frac{m_Sm_E}{R_E}$ \\
    Stabile Umlaufbahn: $\omega_S=\sqrt{\frac{G  m_E}{r^3}}\qquad v=\omega r$\\
    Geostat. Umlaufbahn: $r_g=\sqrt[3]{G  m_E   \frac{T_g^2}{4\pi^2}} \qquad h_g=r_g-R_E$

    TODO: Keplersche Gesetze?
\end{sectionbox}